\section{Hermitian exponential paths}
\label{sec:hermitian_exponential_paths}

\begin{definition}[Exponential path of a matrix]
    \label{def:hermitian_exponential_path}
    \lean{Matrix.hermitianUnitaryPath}
    \leanok
    For a complex square matrix $H$ and $t\in\R$, set $U_H(t)=e^{itH}$.
\end{definition}

\begin{theorem}[Continuous unitary one-parameter group]
    \label{thm:hermitian_exponential_group}
    \lean{Matrix.hermitianUnitaryPath_mul_conjTranspose,
        Matrix.hermitianUnitaryPath_zero,
        Matrix.continuous_hermitianUnitaryPath,
        Matrix.hermitianUnitaryPath_add,
        Matrix.hermitianUnitaryRepresentation}
    \leanok
    \uses{def:hermitian_exponential_path}
    For every complex square matrix $H$, the path $U_H$ is continuous,
    $U_H(0)=\Id$, and $U_H(s+t)=U_H(s)U_H(t)$ for all real $s,t$.
    If $H=H^\dagger$, then $U_H(t)U_H(t)^\dagger=\Id$, so $U_H$
    is a continuous homomorphism from $(\R,+)$ to the unitary group.
\end{theorem}
\begin{proof}\leanok
    Continuity follows from continuity of the matrix exponential.
    The exponentials of $isH$ and $itH$ multiply to $e^{i(s+t)H}$
    because the two exponents commute. For Hermitian $H$,
    $(itH)^\dagger=-itH$, so $U_H(t)^\dagger=e^{-itH}$ and
    $U_H(t)U_H(t)^\dagger=e^0=\Id$.
\end{proof}

\begin{theorem}[Nonscalar exponentials near the identity]
    \label{thm:exponential_nonscalar_near_zero}
    \lean{Matrix.eventually_hermitianUnitaryPath_nonscalar}
    \leanok
    \uses{def:hermitian_exponential_path}
    Let $K$ be a complex square matrix with $K\notin\C\Id$.
    There exists $\varepsilon>0$ such that
    $e^{itK}\notin\C\Id$ whenever $0<|t|<\varepsilon$.
    Hermiticity of $K$ is not required.
\end{theorem}
\begin{proof}\leanok
    Since the common commutant of the matrix units consists of scalar
    matrices, there is a matrix unit $B$ with $BK-KB\ne0$.
    Define $f(t)=Be^{itK}-e^{itK}B$. Then $f(0)=0$ and
    $f'(0)=i(BK-KB)\ne0$. Hence $f(t)/t\to f'(0)$ as $t\to0$ through
    nonzero values, and $f(t)\ne0$ throughout a sufficiently small
    punctured neighborhood. A scalar $e^{itK}$ would instead give $f(t)=0$.
\end{proof}

\begin{theorem}[Exponentials of tensor-factor actions]
    \label{thm:exponential_kronecker_sum}
    \lean{Matrix.exp_kronecker_sum, Matrix.exp_kronecker_one,
        Matrix.exp_one_kronecker, Matrix.exp_map_conjugate,
        Matrix.exp_rowCommutator}
    \leanok
    \uses{def:hermitian_exponential_path}
    For complex square matrices $X,Y$, of possibly different sizes,
    \begin{align}
        e^{X\otimes\Id+\Id\otimes Y}&=e^X\otimes e^Y,\notag\\
        e^{X\otimes\Id}&=e^X\otimes\Id,\notag\\
        e^{\Id\otimes Y}&=\Id\otimes e^Y,\notag\\
        e^{\overline X}&=\overline{e^X}.
        \notag
    \end{align}
    In particular, for every complex square matrix $H$ and $t\in\R$,
    \begin{align}
        e^{it(H\otimes\Id-\Id\otimes\overline H)}
        &=U_H(t)\otimes\overline{U_H(t)}.
        \label{eq:semigroup_doubled_exponential}
    \end{align}
\end{theorem}
\begin{proof}\leanok
    The continuous homomorphisms $X\mapsto X\otimes\Id$ and
    $Y\mapsto\Id\otimes Y$ preserve the exponential series.
    The two tensor-factor actions commute, so the exponential addition
    formula gives the first identity. Transposition and adjunction
    preserve the exponential, and their composition is entrywise
    conjugation. Finally, $\overline{itH}=-it\overline H$; applying
    the tensor-sum identity gives (\ref{eq:semigroup_doubled_exponential}).
\end{proof}
